The field of Quantum Machine Learning (QML) has seen rapid growth, but the susceptibility of Quantum Neural Networks (QNNs) to adversarial perturbations remains a critical challenge. As established by Lu et al. \cite{Lu2020}, quantum classifiers are vulnerable to gradient-based attacks, mirroring the vulnerabilities found in classical deep learning. Among these, the Projected Gradient Descent (PGD) algorithm has emerged as a particularly potent white-box attack capable of circumventing simple defenses \cite{Weber2025}. 

While much research has focused on ideal simulators, evaluating these attacks on the Noisy Intermediate-Scale Quantum (NISQ) devices used today is essential. Qiskit's fake backends offer a vital tool for this purpose, providing noise models derived from real system snapshots \cite{Ren2024, QiskitAer}. However, the fidelity of these models and the generalization properties of adversarially trained quantum models under specific noise profiles require further investigation \cite{Du2024}. This project builds upon the foundational work on QML robustness by Mau{\ss}ner and Reers \cite{Maussner2025}, extending their analysis to a broader range of simulated architectures.

The software implementation for this study, including all benchmarking scripts and the PennyLane-Qiskit integration, is available in the project repository \cite{qml_cysec_repo}.

In this paper, we aim to:
\begin{enumerate}
    \item Evaluate the performance and scalability of Qiskit's Fake Backends (V2) for density matrix simulations, focusing on thread parallelization and Aer method selection.
    \item Quantify the runtime scaling behavior from 4 to 8 qubits to establish the feasibility of medium-scale noisy simulation.
    \item Benchmark the effectiveness of the PGD attack algorithm across multiple 4-qubit hardware-mimicking architectures (\texttt{FakeGuadalupeV2}, \texttt{FakeLimaV2}, \texttt{FakeJakartaV2}).
    \item Analyze the impact of varying adversarial retraining ratios on the robustness of QNNs under realistic noise conditions.
\end{enumerate}
